\chapter{Podstawy teoretyczne}
\label{chapt:Podstawy}

\section{Sygnał EEG i jego rejestracja}
TODO: Pochodzenie sygnału, system 10-20, elektrody mokre i suche, parametry sprzętu konsumenckiego.

\section{Rytmy mózgowe i ich interpretacja}
TODO: Pasma delta, theta, alfa, SMR, beta, gamma; reaktywność alfa; wskaźniki pasmowe (np. stosunek beta/(alfa+theta)).

\section{Neurofeedback}
TODO: Zasada działania (warunkowanie instrumentalne), protokoły treningowe, przegląd dowodów skuteczności i ich ograniczeń.

\section{Artefakty w sygnale EEG}
TODO: Mrugnięcia i ruchy oczu (EOG), aktywność mięśniowa (EMG), ruch, zakłócenia sieciowe 50 Hz; metody detekcji i odrzucania.

\section{Przetwarzanie sygnału w czasie rzeczywistym}
TODO: Filtry przyczynowe IIR/FIR, okna z nakładaniem, estymacja widma (Welch), źródła opóźnień w łańcuchu przetwarzania.
